\documentclass[11pt, letterpaper]{article}

\topmargin 0in \textheight 9in \textwidth 6.7in \evensidemargin -.21in \oddsidemargin -.21in \columnsep 0.3in

\usepackage{amsmath}
\usepackage{amsfonts}
\usepackage{amssymb}
\usepackage{graphicx}
\usepackage{array}
\usepackage{booktabs}
\usepackage{enumitem}
\usepackage{hyperref}
\usepackage{xcolor}
\usepackage{natbib}

\hypersetup{
    colorlinks=true,
    linkcolor=blue!70!black,
    urlcolor=blue!70!black,
    citecolor=blue!70!black
}

\setlength{\parindent}{0pt}
\setlength{\parskip}{6pt}

\begin{document}

\pagestyle{empty}

~\vskip1em
{\Large
\noindent \textbf{ECE 657D: Research Project} \\[4pt]
}
{\large
% ===== INSTRUCTIONS =====
% Replace "Your Project Title" with your actual project title.
% Replace "Your Name" and "Your Student ID" below.
\noindent \textbf{Your Project Title Here} \\[4pt]
\noindent Your Name \quad | \quad Your Student ID \\[4pt]
\noindent University of Waterloo \quad | \quad Winter 2026
}

\vskip 1em
\hrule
\vskip 1em

% ===== INSTRUCTIONS =====
% This template follows the adapted Heilmeier Catechism from the project handout.
% Each section corresponds to one of the eight evaluation questions.
%
% FOR THE PROPOSAL (due Feb 27):
%   - Target length: 2 pages (excluding references)
%   - Some sections may be preliminary; that is expected.
%
% FOR THE FINAL SUBMISSION (due Apr 2):
%   - Target length: 4 pages (excluding references)
%   - All sections should be complete, with results and analysis.
%   - Add a "Results" section after Q8 to present your findings.
%
% Delete all instructional comments before submitting.

\section*{Q1: Research Question}
% What is your research question? State it in plain language.
% A reader with no background in your specific topic should understand
% what you are investigating.

[Your research question here. For example: ``Can a lightweight vision transformer match the performance of a ResNet-50 on medical image classification when trained with limited labeled data?'']

\section*{Q2: Prior Work and Limitations}
% Summarize the most relevant existing approaches. What has been tried?
% What are the key limitations that motivate your work? Cite specific papers.

[Describe 3--5 relevant prior works and their limitations. For example: \citet{he2016deep} introduced residual connections that enabled training of very deep networks, but their approach requires large labeled datasets. \citet{dosovitskiy2021image} showed that transformers can achieve competitive results on image classification, but require extensive pre-training data.]

\section*{Q3: Novelty of Approach}
% What is new or different about your method?
% This could be a novel combination of techniques, an application to a new domain,
% or a systematic comparison that has not been done before.

[Describe what is novel about your approach here.]

\section*{Q4: Impact and Motivation}
% Who would benefit from your results? Why should we care?
% Connect your work to broader goals in deep learning or its applications.

[Explain why this matters and who would benefit.]

\section*{Q5: Risks and Assumptions}
% What are the main risks? What could go wrong?
% What assumptions about data, compute, or method might not hold?

[Identify key risks and assumptions here.]

\section*{Q6: Resources}
% What compute and data resources do you need?
% What datasets will you use? Where will you get compute?

[Describe your resource plan. For example: ``I will use the CIFAR-10 dataset \citep{krizhevsky2009learning} and train on Google Colab free tier (T4 GPU). Training is estimated to take approximately 4 hours per experiment.'']

\section*{Q7: Weekly Plan}
% Provide a week-by-week plan for the ~5 weeks between proposal and final submission.

\begin{center}
\renewcommand{\arraystretch}{1.3}
\begin{tabular}{cl}
\toprule
\textbf{Week} & \textbf{Plan} \\
\midrule
1 (Mar 3--7) & [e.g., Finalize dataset, implement baseline model] \\
2 (Mar 10--14) & [e.g., Implement proposed method] \\
3 (Mar 17--21) & [e.g., Run experiments, tune hyperparameters] \\
4 (Mar 24--28) & [e.g., Additional experiments, ablation studies] \\
5 (Mar 31--Apr 2) & [e.g., Write up results, finalize submission] \\
\bottomrule
\end{tabular}
\end{center}

\section*{Q8: Milestones and Success Criteria}
% How will you measure success? What baselines will you compare against?
% Define quantitative and/or qualitative metrics.

[Define your success criteria here. For example: ``Success is measured by test accuracy on CIFAR-10. The baseline is a ResNet-18 trained from scratch. I will consider the project successful if the proposed method achieves within 2\% of the baseline accuracy with 50\% fewer parameters.'']

% ===== FOR THE FINAL SUBMISSION ONLY =====
% Uncomment and fill in the sections below for your final write-up.

% \section*{Results}
% [Present your experimental results here. Include tables and/or figures.]

% \section*{Discussion}
% [Discuss your results. What worked? What didn't? What did you learn?]

% \section*{Acknowledgments}
% [Acknowledge any AI tools used, discussions with classmates, etc.]

\bibliographystyle{plainnat}
\bibliography{references}

\end{document}
